\documentclass[a4paper,11pt]{article}

% ─── Geometry ────────────────────────────────────────────────────────────────
\usepackage{geometry}
\geometry{margin=1in, headheight=15pt}

% ─── Mathematics ─────────────────────────────────────────────────────────────
\usepackage{amsmath}
\usepackage{amssymb}
\usepackage{mathtools}

% ─── Graphics / Colours ──────────────────────────────────────────────────────
\usepackage{graphicx}
\usepackage{xcolor}
\definecolor{blue3}{RGB}{0,0,180}
\definecolor{green3}{RGB}{0,120,0}
\definecolor{orange3}{RGB}{200,100,0}
\definecolor{red3}{RGB}{180,0,0}
\definecolor{grey3}{RGB}{80,80,80}

% ─── Units ───────────────────────────────────────────────────────────────────
\usepackage{siunitx}
\sisetup{detect-all}

% ─── Links (hyperref BEFORE cleveref) ────────────────────────────────────────
\usepackage[colorlinks=true,linkcolor=blue3,citecolor=blue3,urlcolor=blue3]{hyperref}
\usepackage{cleveref}

% ─── Boxes ───────────────────────────────────────────────────────────────────
\usepackage{tcolorbox}
\tcbuselibrary{skins,breakable}

\newtcolorbox{pointsbox}[1][]{
    colback=grey3!8,
    colframe=grey3,
    fonttitle=\bfseries\small,
    #1
}

% ─── Tables / Lists ──────────────────────────────────────────────────────────
\usepackage{booktabs}
\usepackage{enumitem}

% ─── Notation commands ───────────────────────────────────────────────────────
\newcommand{\ktilde}{\tilde{k}}
\newcommand{\ytilde}{\tilde{y}}
\newcommand{\ctilde}{\tilde{c}}
\newcommand{\dd}{\mathrm{d}}
\newcommand{\bgp}{\mathrm{BGP}}

% ─── Header / Footer ─────────────────────────────────────────────────────────
\usepackage{fancyhdr}
\pagestyle{fancy}
\fancyhf{}
\lhead{\textbf{14E022 Macroeconomics I} \quad Practice Exam 5}
\rhead{BSE 2025--26}
\cfoot{\thepage}

% ═════════════════════════════════════════════════════════════════════════════
\begin{document}
% ═════════════════════════════════════════════════════════════════════════════

\begin{center}
    {\LARGE \textbf{14E022 Macroeconomics I}}\\[6pt]
    {\Large \textbf{Practice Exam 5}}\\[8pt]
    {\large BSE --- Academic Year 2025--26}
\end{center}

\vspace{8pt}
\noindent\textit{Instructions: the exam consists of 3 problems and a set of short questions
totalling \textbf{100 points}. You have \textbf{2 hours}. Show all derivations clearly.
For short questions, concise and precise answers score highest.}
\vspace{4pt}

\noindent\rule{\textwidth}{0.4pt}

% ─────────────────────────────────────────────────────────────────────────────
\section{Question 1 --- Ramsey Model with Labour Income Tax \hfill (25 points)}
% ─────────────────────────────────────────────────────────────────────────────

Consider a standard Ramsey--Cass--Koopmans economy with no population growth ($n=0$) and
no exogenous TFP growth. The representative household has preferences
\[
    \int_0^\infty e^{-\rho t}\frac{c_t^{1-\sigma}-1}{1-\sigma}\,\dd t,
    \qquad \sigma > 0,\;\rho > 0.
\]
The household inelastically supplies one unit of labour. The government levies a
proportional tax $\tau \in (0,1)$ on \emph{labour income} and returns all revenue as a
lump-sum transfer $T_t = \tau w_t$ to the household. There is no tax on capital income.
Firms operate with a neoclassical production function $Y_t = F(K_t, L_t)$ satisfying CRS
and positive, diminishing marginal products. Write $f(k) = F(k,1)$ and
$r_t = f'(k_t) - \delta$.

\begin{enumerate}[label=\textbf{\arabic*.}, leftmargin=*, itemsep=10pt]

    \item \textbf{(6 points)} Write the household's flow budget constraint, accounting for
    the lump-sum transfer. Set up the current-value Hamiltonian and derive the Euler
    equation for consumption. Show that the tax $\tau$ does \emph{not} appear in the Euler
    equation, and explain intuitively why.

    \item \textbf{(7 points)} Characterise the balanced growth path (BGP). Write the two
    conditions that define the steady state $(k^*, c^*)$: the modified Golden Rule for
    capital and market clearing for consumption. Compare $k^*$ to the no-tax Ramsey steady
    state and to the Golden Rule capital level $k^{GR}$. Is the labour income tax
    distortionary in this model? Justify your answer.

    \item \textbf{(6 points)} Suppose the economy is at the no-tax Ramsey steady state when,
    unexpectedly at $t=0$, the government announces and immediately implements $\tau > 0$.
    Describe the transition dynamics. What happens to $c$ and $k$ on impact and along the
    transition path? Use the phase diagram to support your answer.

    \item \textbf{(6 points)} Now suppose the government uses the tax revenue to finance a
    public good $G_t = \tau w_t$ that does not affect production or household utility
    directly (instead of rebating as a lump-sum transfer). How do your answers to parts
    1--3 change? Is the labour income tax now distortionary? How does the steady state
    differ?

\end{enumerate}

\clearpage

% ─────────────────────────────────────────────────────────────────────────────
\section{Question 2 --- Aghion--Howitt Quality Ladders \hfill (25 points)}
% ─────────────────────────────────────────────────────────────────────────────

Consider the Aghion--Howitt (1992) model of Schumpeterian growth. There is a single
intermediate good whose current quality level is $q_t$. Final output is
\[
    Y_t = A\,q_t\,x_t^{\alpha}, \qquad \alpha \in (0,1),
\]
where $x_t$ is the quantity of the intermediate good and $A > 0$ is a constant.
The intermediate good is produced one-for-one from labour. A successful innovation arrives
at aggregate rate $\iota$ (the creative destruction rate chosen by entrant researchers).
Each unit of research labour generates a Poisson innovation at rate $\lambda > 0$.
The new quality leader sets quality $\lambda q_t$ (multiplicative improvement,
$\lambda > 1$). The interest rate is $r$. Labour market clears:
\[
    x_t + L_{R,t} = L,
\]
where $L_{R,t}$ is research labour and $L$ is fixed total labour supply.

\begin{enumerate}[label=\textbf{\arabic*.}, leftmargin=*, itemsep=10pt]

    \item \textbf{(5 points)} Derive the monopoly price $p^*$ and monopoly profits $\pi$
    for the quality leader. Express $\pi$ in terms of $\alpha$, $\lambda$, $A$, $q_t$, and
    the wage $w$.

    \item \textbf{(7 points)} Write the no-arbitrage Bellman equation for the value $V$ of
    being the current quality leader. Solving on the BGP (where $\dot{V}=0$ and $r$, $\iota$,
    $\pi$ are constant), find the equilibrium value $V^*$. Write the free-entry condition
    for researchers. Derive the equilibrium creative destruction rate $\iota^*$ as a function
    of the model parameters.

    \item \textbf{(5 points)} Identify and precisely define the two key externalities in
    this model: the \textbf{business-stealing externality} and the
    \textbf{appropriability (consumer-surplus) externality}. State the direction of each
    distortion (does it push $\iota$ above or below the social optimum?) and explain the
    economic intuition.

    \item \textbf{(4 points)} Under what condition on the model parameters does equilibrium
    R\&D ($\iota^*$) exceed the socially optimal level? Write the condition and explain
    which externality dominates when this holds.

    \item \textbf{(4 points)} Propose a policy mix that could implement the social optimum.
    Explain why a single instrument is generally insufficient.

\end{enumerate}

\clearpage

% ─────────────────────────────────────────────────────────────────────────────
\section{Question 3 --- Development Accounting with Heterogeneous Human Capital
\hfill (25 points)}
% ─────────────────────────────────────────────────────────────────────────────

Output of country $i$ is
\[
    Y_i = K_i^{\alpha}(h_i L_i)^{1-\alpha} A_i, \qquad \alpha = 1/3.
\]
Two countries are observed:

\vspace{6pt}
\begin{center}
\begin{tabular}{lccc}
\toprule
Country & $Y_i/L_i$ (normalised) & $K_i/Y_i$ & $h_i$ \\
\midrule
US (benchmark) & 1.00 & 3.00 & 1.00 \\
Country E       & 0.10 & 2.00 & 0.55 \\
\bottomrule
\end{tabular}
\end{center}
\vspace{6pt}

\begin{enumerate}[label=\textbf{\arabic*.}, leftmargin=*, itemsep=10pt]

    \item \textbf{(6 points)} Starting from the production function, derive the development
    accounting decomposition expressing $y_i = Y_i/L_i$ as a product of a
    capital-intensity factor, a human-capital factor, and TFP. Be explicit about the
    algebra.

    \item \textbf{(8 points)} Using the data, compute the fraction of the income gap (in
    logs) between Country E and the US explained by: (a) physical capital intensity,
    (b) human capital, (c) TFP. Show all calculations.

    \item \textbf{(6 points)} A referee argues: ``The large TFP residual does not reflect
    true technology differences --- it simply captures resource misallocation.'' Explain
    this critique carefully with reference to the Hsieh--Klenow (2009) framework. What
    does TFPR dispersion reveal, and how does it affect measured aggregate TFP?

    \item \textbf{(5 points)} Given that TFP explains the bulk of income gaps, and given
    the misallocation critique, what are the implications for development policy? Discuss
    two specific policy recommendations and their limitations.

\end{enumerate}

\clearpage

% ─────────────────────────────────────────────────────────────────────────────
\section{Short Questions \hfill (25 points)}
% ─────────────────────────────────────────────────────────────────────────────

\noindent\textit{For these questions, short and sharp answers will receive the highest marks.
Aim for 3--6 sentences per question unless otherwise indicated.}

\vspace{8pt}

\begin{enumerate}[label=\textbf{\arabic*.}, leftmargin=*, itemsep=14pt]

    % ── SQ1 ──────────────────────────────────────────────────────────────────
    \item \textbf{(8 points) Kaldor Facts and Balanced Growth.}

    State \emph{three} Kaldor stylised facts about long-run growth. Which families of
    production functions are consistent with balanced growth (constant growth rates and
    constant factor shares) and why? Why does the Cobb-Douglas production function play a
    privileged role among these?

    % ── SQ2 ──────────────────────────────────────────────────────────────────
    \item \textbf{(9 points) TVC and the No-Ponzi Condition.}

    Define the Transversality Condition (TVC) precisely in the Ramsey model. Define the
    No-Ponzi condition on household debt. Show formally how the TVC rules out
    over-accumulation (dynamic inefficiency) in the Ramsey model. Contrast this with the
    Solow model, where dynamic inefficiency is possible.

    % ── SQ3 ──────────────────────────────────────────────────────────────────
    \item \textbf{(8 points) Baumol Cost Disease and Aggregate Productivity.}

    What is Baumol's cost disease? Explain the mechanism: why does the stagnant sector's
    relative price rise over time, and what drives rising employment in that sector?
    Discuss whether rising service employment necessarily implies a fall in aggregate
    productivity growth, and under what preference assumption the employment share of
    services stays constant over time.

\end{enumerate}

\vspace{20pt}
\noindent\rule{\textwidth}{0.4pt}

\begin{center}
    \textit{End of Practice Exam 5. Good luck!}
\end{center}

% ─────────────────────────────────────────────────────────────────────────────
\clearpage
\appendix
\section*{Appendix: Key Formulas Permitted in Exam}
% ─────────────────────────────────────────────────────────────────────────────

The following results may be used without re-derivation unless the question explicitly
asks you to derive them.

\subsection*{A.1 \enspace Log-differentiation}
If $Z_t = X_t^a Y_t^b$, then $\gamma_Z = a\gamma_X + b\gamma_Y$.

\subsection*{A.2 \enspace Ramsey / Euler Equation (standard, $n=0$)}
\[
    \frac{\dot{c}}{c} = \frac{1}{\sigma}\bigl(r - \rho\bigr),
    \qquad r = f'(k) - \delta.
\]

\subsection*{A.3 \enspace Hamiltonian first-order conditions (current-value)}
Control $u$, state $x$, costate $\mu$, discount $\rho$:
\[
    \frac{\partial H}{\partial u} = 0, \qquad
    \dot{\mu} = \rho\mu - \frac{\partial H}{\partial x}.
\]

\subsection*{A.4 \enspace Cobb-Douglas factor payments}
$rK = \alpha Y$, $wL = (1-\alpha)Y$, so $r = \alpha Y/K$ and $w = (1-\alpha)Y/L$.

\subsection*{A.5 \enspace Development accounting identity}
With $y_i = Y_i/L_i$, $\alpha = 1/3$, and capital-output ratio $\kappa_i = K_i/Y_i$:
\[
    \frac{y_i}{y_{US}} = \left(\frac{\kappa_i}{\kappa_{US}}\right)^{\!\alpha/(1-\alpha)}
    \cdot \frac{h_i}{h_{US}}
    \cdot \frac{A_i}{A_{US}}.
\]

\subsection*{A.6 \enspace Monopoly markup and profits (Aghion--Howitt)}
With intermediate production one-for-one from labour at wage $w$ and final-good demand
$x = (\alpha A q / p)^{1/(1-\alpha)}$:
\[
    p^* = \frac{w}{\alpha}, \qquad
    \pi = \frac{1-\alpha}{\alpha}\,w\,x^*.
\]

\subsection*{A.7 \enspace Aghion--Howitt Bellman on BGP ($\dot{V}=0$)}
\[
    V^* = \frac{\pi}{r + \iota}.
\]

\end{document}
